% !TeX root = ../../Book3.tex
\section{Euler \texorpdfstring{\(\Gamma\)}{\AnalysisPdfGamma} 函数}
\label{b3:sec:2101}

阶乘满足乘法递推。若要把整数变量换成实变量或复变量，需要先找到一个在整数处给出阶乘、同时又能对参数作分析运算的表达。Euler 积分提供了这样的函数，但递推和整数点取值并不足以唯一确定它。本节先从积分得到全纯性、递推和半整数值，再用正实轴上的对数凸性建立唯一性与平移比值估计。

\subsection{Euler 积分}
\label{b3:subsec:210101}

积分 \(\int_0^\infty e^{-t}t^n\,\mathrm{d}t\) 中，分部积分会把幂次降低一，并产生因子 \(n\)。为了使递推写成变量本身乘以函数，定义时把幂次取为 \(z-1\)。

\Needspace{6\baselineskip}
\begin{definition}\label{b3:def:21-gamma}
对 \(z\in\mathbb C\)、\(\operatorname{Re}z>0\)，定义\term[gamma-hanshu]{Gamma 函数}[gamma function]
\[
 \Gamma(z)\coloneqq\int_0^\infty e^{-t}t^{z-1}\,\mathrm{d}t,
 \quad t^{z-1}=\exp\bigl((z-1)\log t\bigr),\quad t>0,
\]
其中 \(\log t\) 是正实轴上的实对数。
\end{definition}
\symidx{\(\Gamma(z)\)：Gamma 函数}
积分变量只取正实数，所以复幂在这里没有绕过原点选择分支的问题。若 \(z=\sigma+\mathrm{i}y\)，则 \(|t^{z-1}|=t^{\sigma-1}\)，虚部只产生模为一的振荡因子。

\ProseParagraphBreak

两端的收敛机制不同。在 \((0,1)\) 上，\(e^{-1}\le e^{-t}\le1\)，所以绝对可积恰好要求 \(\sigma>0\)。在 \([1,\infty)\) 上，指数衰减压过任意固定实次幂，对所有实数 \(\sigma\) 都可积。因此 \(\operatorname{Re}z>0\) 是这个积分的绝对收敛域。要在积分号下对 \(z\) 求导，还须控制每次求导带来的 \(\log t\)。

\Needspace{6\baselineskip}
\begin{proposition}[全纯性与参数导数]\label{b3:prop:21-euler-derivatives}
Euler 积分定义的 \(\Gamma\) 在半平面 \(\operatorname{Re}z>0\) 全纯。对每个非负整数 \(k\)，有
\[
 \Gamma^{(k)}(z)
 =\int_0^\infty e^{-t}t^{z-1}(\log t)^k\,\mathrm{d}t,
 \quad \operatorname{Re}z>0,
\]
且这些积分在每个紧参数集上绝对一致收敛。此外，\(\Gamma(x)>0\) 对 \(x>0\) 成立，并且
\[
 |\Gamma(\sigma+\mathrm{i}y)|\le\Gamma(\sigma),
 \quad \sigma>0,\quad y\in\mathbb R.
\]
\end{proposition}
\begin{proof}
固定紧集 \(K\subset\{\operatorname{Re}z>0\}\)，选 \(0<a\le\operatorname{Re}z\le b\) 对所有 \(z\in K\) 成立。在 \(0<t<1\) 上，
\[
 |e^{-t}t^{z-1}(\log t)^k|
 \le t^{a-1}|\log t|^k.
\]
代换 \(u=-\log t\)，再作有限次分部积分，得到
\[
 \int_0^1t^{a-1}|\log t|^k\,\mathrm{d}t
 =\int_0^\infty e^{-au}u^k\,\mathrm{d}u
 =\frac{k!}{a^{k+1}}<\infty.
\]
在 \(t\ge1\) 上，则有
\[
 |e^{-t}t^{z-1}(\log t)^k|
 \le e^{-t}t^{b-1}(\log t)^k
 \le C_{b,k}e^{-t/2}.
\]
最后一个常数存在，因为 \(t^{b-1}(\log t)^k e^{-t/2}\) 在 \([1,\infty)\) 有界。这些控制函数不依赖 \(z\in K\)。

截断积分 \(G_N(z)=\int_{1/N}^{N}e^{-t}t^{z-1}\,\mathrm{d}t\) 是整函数，在有限积分区间可以逐次对参数求导。上述共同控制表明，\(G_N\) 及其各个固定阶导数的积分表达在 \(K\) 上一致收敛。由\cref{b3:prop:04-derivative-convergence}，极限全纯，导数正是陈述中的积分。正实参数时，被积函数处处为正，故 \(\Gamma(x)>0\)。最后，对 Euler 积分使用三角不等式即得竖线上的模估计。
\end{proof}
这里的一致性针对紧参数集；当实部趋近零时，共同控制中的 \(a\) 也会趋近零，不能保留同一个可积上界。右半平面上的竖线估计只给出有界性，尚未给出随 \(|y|\) 增大的衰减速度。

\subsection{函数方程}
\label{b3:subsec:210102}

积分定义使变量平移一步恰好对应一次分部积分。端点项是否消失，仍由原点的幂次和无穷远的指数分别控制。

\Needspace{6\baselineskip}
\begin{proposition}[递推与归一化]\label{b3:prop:21-recurrence}
对所有 \(z\in\mathbb C\)、\(\operatorname{Re}z>0\)，有
\[
 \Gamma(z+1)=z\Gamma(z),\quad \Gamma(1)=1.
\]
因此，对每个非负整数 \(n\)，\(\Gamma(n+1)=n!\)。
\end{proposition}
\begin{proof}
在 \([\varepsilon,R]\) 上分部积分，得到
\[
 \int_\varepsilon^R e^{-t}t^z\,\mathrm{d}t
 =\bigl[-e^{-t}t^z\bigr]_\varepsilon^R
  +z\int_\varepsilon^R e^{-t}t^{z-1}\,\mathrm{d}t.
\]
设 \(\sigma=\operatorname{Re}z>0\)。无穷远端的边界项模为 \(e^{-R}R^\sigma\to0\)，原点端的模为 \(e^{-\varepsilon}\varepsilon^\sigma\to0\)。两个积分都绝对收敛，所以可分别令 \(R\to\infty\)、\(\varepsilon\downarrow0\)，得到递推式。直接计算 \(\int_0^\infty e^{-t}\,\mathrm{d}t=1\) 给出归一化，再作整数次递推得到阶乘。
\end{proof}
归一化选定了乘法常数，却没有排除一周期因子。例如，固定 \(0<\eta<1\)，函数
\[
 F(z)=\Gamma(z)\bigl(1+\eta\sin(2\pi z)\bigr)
\]
在右半平面全纯，满足同一递推、同一整数点取值，而且在正实轴仍为正。它与 \(\Gamma\) 不同，因此即使同时要求全纯性和正实轴上的正性，递推仍不能给出唯一性。

\ProseParagraphBreak

递推还把半整数值归约到一个起始值。这个值由 Gaussian 积分确定，适合在讨论渐近公式之前独立算出。

\Needspace{6\baselineskip}
\begin{proposition}[Gaussian 积分与半整数值]\label{b3:prop:21-half-integers}
对每个非负整数 \(n\)，有
\[
 \Gamma\left(\frac12\right)=\sqrt\pi,
 \quad
 \Gamma\left(n+\frac12\right)
 =\frac{(2n)!}{4^n n!}\sqrt\pi.
\]
\end{proposition}
\begin{proof}
令 \(I=\int_{\mathbb R}e^{-u^2}\,\mathrm{d}u\)。这个积分有限且为正：在 \(|u|\ge1\) 上可由 \(e^{-|u|}\) 控制，在有限区间被积函数有界。被积函数非负，Tonelli 定理给
\[
 I^2=\int_{\mathbb R^2}e^{-(u^2+v^2)}\,\mathrm{d}u\,\mathrm{d}v.
\]
在圆环上使用极坐标换元，Jacobi 因子为 \(r\)，再让内半径趋零、外半径趋于无穷。非负性允许用单调收敛取极限，于是
\[
 I^2=\int_0^{2\pi}\int_0^\infty e^{-r^2}r\,\mathrm{d}r\,\mathrm{d}\theta
 =2\pi\cdot\frac12=\pi.
\]
由 \(I>0\) 确定 \(I=\sqrt\pi\)。在 Euler 积分中令 \(t=u^2\)，得到
\[
 \Gamma\left(\frac12\right)
 =2\int_0^\infty e^{-u^2}\,\mathrm{d}u=I.
\]
最后，重复使用递推式，
\[
 \Gamma\left(n+\frac12\right)
 =\left(\prod_{j=0}^{n-1}\left(j+\frac12\right)\right)\sqrt\pi
 =\frac{1\cdot3\cdots(2n-1)}{2^n}\sqrt\pi
 =\frac{(2n)!}{4^n n!}\sqrt\pi.
\]
\(n=0\) 时有限乘积为空，公式仍成立。
\end{proof}
同一换元给出 Gaussian 偶数矩
\[
 \int_{\mathbb R}u^{2n}e^{-u^2}\,\mathrm{d}u
 =\Gamma\left(n+\frac12\right),\quad n\ge0.
\]
奇数矩因被积函数为奇函数而为零；这些积分均绝对收敛。阶乘和半整数值由此分别对应指数权重与 Gaussian 权重下的整数矩。

\subsection{对数凸性}
\label{b3:subsec:210103}

Euler 积分中的参数出现在指数 \(t^{x-1}\) 上，因此参数的凸组合会变成两个正被积函数的几何平均。Hölder 不等式正好比较这种几何平均的积分。

\Needspace{6\baselineskip}
\begin{proposition}[严格对数凸性]\label{b3:prop:21-logconvex}
Euler 积分定义的实函数 \(x\mapsto\log\Gamma(x)\) 在 \((0,\infty)\) 上严格凸。即对 \(x,y>0\)、\(x\ne y\) 和 \(0<\lambda<1\)，
\[
 \Gamma(\lambda x+(1-\lambda)y)
 <\Gamma(x)^\lambda\Gamma(y)^{1-\lambda}.
\]
\end{proposition}
\begin{proof}
令 \(A(t)=e^{-t}t^{x-1}\)、\(B(t)=e^{-t}t^{y-1}\)。二者严格为正且可积，并且
\[
 e^{-t}t^{\lambda x+(1-\lambda)y-1}
 =A(t)^\lambda B(t)^{1-\lambda}.
\]
对右边使用指数 \(1/\lambda\) 与 \(1/(1-\lambda)\) 的 Hölder 不等式，得
\[
 \int_0^\infty A^\lambda B^{1-\lambda}\,\mathrm{d}t
 \le\left(\int_0^\infty A\,\mathrm{d}t\right)^\lambda
      \left(\int_0^\infty B\,\mathrm{d}t\right)^{1-\lambda}.
\]
在两个非零正函数之间，等号要求 \(A\) 与 \(B\) 几乎处处成比例，即 \(t^{x-y}\) 几乎处处为常数。但 \(x\ne y\) 时该函数在 \((0,\infty)\) 严格单调，不可能满足这个条件。因此不等式严格，取对数即得结论。
\end{proof}
还可以把严格性读成一个二阶导数公式。固定 \(x>0\)，以
\[
 w_x(t)=\frac{e^{-t}t^{x-1}}{\Gamma(x)},\quad t>0,
\]
作为概率密度，并令 \(m_x=\int_0^\infty(\log t)w_x(t)\,\mathrm{d}t\)。参数导数公式保证这些对数矩有限，且
\[
 \frac{\mathrm{d}^2}{\mathrm{d}x^2}\log\Gamma(x)
 =\frac{\Gamma''(x)}{\Gamma(x)}
  -\left(\frac{\Gamma'(x)}{\Gamma(x)}\right)^2
 =\int_0^\infty(\log t-m_x)^2w_x(t)\,\mathrm{d}t>0.
\]
最后一项是 \(\log t\) 的方差；密度在整个正半轴严格为正，而 \(\log t\) 不恒定，所以方差严格为正。归一化的分母 \(\Gamma(x)\) 是这一步所必需的，单独观察 \(\Gamma''(x)>0\) 只能得到普通凸性。

\subsection{Bohr–Mollerup 定理}
\label{b3:subsec:210104}

递推固定了整数点的对数差值，而凸性控制整数之间的割线斜率。把一个未知的 \(F(x)\) 向右递推到 \(F(n+x)\)，再利用相邻整数之间的斜率夹逼，就能让未知量落入相对宽度趋于零的区间。这给出\term[bohr-mollerup]{Bohr--Mollerup 定理}[Bohr--Mollerup theorem]。

\Needspace{6\baselineskip}
\begin{theorem}[Bohr--Mollerup]\label{b3:thm:21-bohr-mollerup}
设 \(F:(0,\infty)\to(0,\infty)\) 满足
\[
 F(1)=1,\quad F(x+1)=xF(x)\quad(x>0),
\]
并且 \(x\mapsto\log F(x)\) 凸。则对每个 \(x>0\)，\(F(x)=\Gamma(x)\)。
\end{theorem}
\begin{proof}
令 \(h=\log F\)。递推与归一化先给 \(F(n)=(n-1)!\)，因而对整数 \(n\ge2\)，
\[
 h(n)-h(n-1)=\log(n-1),\quad
 h(n+1)-h(n)=\log n.
\]
固定 \(0<x\le1\)。凸函数的割线斜率随区间向右移动而非减，所以
\[
 h(n)-h(n-1)
 \le\frac{h(n+x)-h(n)}x
 \le h(n+1)-h(n).
\]
右侧亦可由 \(n+x=(1-x)n+x(n+1)\) 的凸性直接得到。于是
\[
 x\log(n-1)\le h(n+x)-h(n)\le x\log n.
\]
指数化后，再用 \(F(n+x)=F(x)x(x+1)\cdots(x+n-1)\)，可得
\[
 A_n(x)\le F(x)\le B_n(x),
\]
其中
\[
 A_n(x)=\frac{(n-1)!(n-1)^x}{x(x+1)\cdots(x+n-1)},\quad
 B_n(x)=\frac{(n-1)!n^x}{x(x+1)\cdots(x+n-1)}.
\]
这两个端点都严格为正，且 \(B_n(x)/A_n(x)=(n/(n-1))^x\to1\)。

已经由积分证明的 \(\Gamma\) 同样满足归一化、递推和对数凸性，所以也落在同一个区间 \([A_n(x),B_n(x)]\) 中。两次夹逼相除，得到
\[
 \left(\frac{n-1}{n}\right)^x
 \le\frac{F(x)}{\Gamma(x)}
 \le\left(\frac n{n-1}\right)^x.
\]
令 \(n\to\infty\)，即得 \(F(x)=\Gamma(x)\) 于 \((0,1]\)。任意正实数可写成 \(x+m\)，其中 \(0<x\le1\)、\(m\ge0\) 为整数，再用有限次递推便扩展到全部正实轴。
\end{proof}
证明只要求 \(\log F\) 凸，不预先要求 \(F\) 可微。若另外给定一个在右半平面全纯的候选函数，且它在正实轴满足定理的条件，则先由本定理识别实轴上的值，再由恒等定理识别整个右半平面内的函数。实变量唯一性与复变量唯一性分别由这两步承担。

\ProseParagraphBreak

对数凸性也能给出精确到相对误差的平移比值。这里的自变量始终沿正实轴趋于无穷，平移参数则保持固定。

\Needspace{6\baselineskip}
\begin{proposition}[固定实平移的比值]\label{b3:prop:21-shift-ratio}
设 \(x>0\)、\(0\le a\le1\)。则
\[
 \left(\frac{x}{x+a}\right)^{1-a}
 \le\frac{\Gamma(x+a)}{x^a\Gamma(x)}\le1.
\]
特别地，对任意固定的 \(a,b\in\mathbb R\)，当 \(x\to+\infty\) 时有
\[
 \frac{\Gamma(x+a)}{\Gamma(x+b)}
 =x^{a-b}\bigl(1+O_{a,b}(x^{-1})\bigr).
\]
后一公式在所有参数 \(x+a,x+b\) 均为正的充分大半轴上使用。
\end{proposition}
\begin{proof}
\(a=0,1\) 时由归一化的比值和递推直接成立。对 \(0<a<1\)，将 \(x+a=(1-a)x+a(x+1)\) 代入对数凸性，得到
\[
 \Gamma(x+a)\le\Gamma(x)^{1-a}\Gamma(x+1)^a
 =x^a\Gamma(x).
\]
另一方面，\(x+1=a(x+a)+(1-a)(x+a+1)\)，故
\[
 \Gamma(x+1)
 \le\Gamma(x+a)^a\Gamma(x+a+1)^{1-a}
 =(x+a)^{1-a}\Gamma(x+a).
\]
使用 \(\Gamma(x+1)=x\Gamma(x)\) 并整理，便得下界。固定 \(a\in[0,1]\) 时，下界为 \((1+a/x)^{-(1-a)}=1+O_a(x^{-1})\)，所以
\[
 \frac{\Gamma(x+a)}{\Gamma(x)}=x^a\bigl(1+O_a(x^{-1})\bigr).
\]

任意固定实数 \(c\) 可写成 \(c=k+\theta\)，其中 \(k\in\mathbb Z\)、\(0\le\theta<1\)。若 \(k\ge0\)，有限次递推给
\[
 \Gamma(x+c)=\Gamma(x+\theta)
             \prod_{j=0}^{k-1}(x+\theta+j).
\]
若 \(k=-q<0\)，则
\[
 \Gamma(x+c)=\frac{\Gamma(x+\theta)}
 {\prod_{j=0}^{q-1}(x+c+j)}.
\]
对充分大的 \(x\)，这些因子均为正；每个因子提出一个 \(x\) 后，有限乘积或其倒数均为 \(1+O_c(x^{-1})\)。结合已经证明的 \(\theta\) 情形，得到 \(\Gamma(x+c)/\Gamma(x)=x^c(1+O_c(x^{-1}))\)。分别取 \(c=a,b\)，再相除，即得一般比值。
\end{proof}
误差常数依赖固定的平移参数；这条结论不能直接用于随 \(x\) 增大的平移量。它也没有给出复变量在任意方向趋于无穷的比值渐近。

\ProseParagraphBreak

半整数值与比值估计合起来，已经能够确定中心二项式系数的主项。对正整数 \(n\)，
\[
 \binom{2n}{n}
 =\frac{4^n}{\sqrt\pi}
   \frac{\Gamma(n+1/2)}{\Gamma(n+1)}
 =\frac{4^n}{\sqrt{\pi n}}\bigl(1+O(n^{-1})\bigr).
\]
这里常数 \(\sqrt\pi\) 来自 Gaussian 积分，比值的幂次来自对数凸性与递推。两种信息共同确定了这个离散系数的渐近尺度。
